\chapter{Qualité logicielle et limites de l'étude}\label{chap:qualite}
\section{Séparer les responsabilités du programme}
Le noyau représente les instances, les paramètres, les solutions et leurs
règles de validité. Les modules de méthodes construisent ou améliorent des
solutions. Une couche d'exécution commune mesure les calculs et produit un
résultat structuré. Les campagnes orchestrent les cas, tandis que les modules
d'artefacts et de visualisation enregistrent les données et dessinent les figures.

\begin{table}[htbp]\small
\begin{tabularx}{\textwidth}{lL}\toprule
Responsabilité & Modules principaux \\\midrule
Données et modèle & \code{instance}, \code{io}, \code{distances}, \code{solution} \\
Validité & \code{validation} \\
Résolution & \code{heuristics}, \code{metaheuristics}, \code{exact} \\
Exécution et campagnes & \code{execution}, \code{experiments}, \code{benchmark} \\
Résultats et figures & \code{artifacts}, \code{visualization}, \code{experiment_plots} \\
Interface & \code{cli}, point d'entrée \code{ring-star} et Makefile \\
\bottomrule\end{tabularx}
\caption{Organisation des responsabilités dans le projet.}
\end{table}

Cette séparation limite les divergences entre une résolution individuelle
et une campagne : les deux empruntent le même chemin d'exécution. Elle
évite également de reconstruire une solution approximative à partir d'une
simple liste de stations lorsqu'un ordre de cycle sauvegardé existe déjà.

\section{Des résultats autonomes et consultables}
Une exécution sauvegarde les coordonnées, les paramètres, les stations,
le cycle, les affectations, les coûts et les informations de statut.
Les recherches ajoutent la solution initiale et une trace. Le journal CBC
est conservé lorsque cette méthode est utilisée. La campagne fournit un
CSV détaillé, une synthèse statistique et un résumé de complétion.

Les dossiers utilisent un identifiant unique pour éviter les écrasements.
L'index global oriente vers les résolutions individuelles et les campagnes,
puis les résumés guident vers les données et figures pertinentes. Les
figures peuvent être régénérées depuis les solutions sauvegardées sans
disposer du fichier TSPLIB initial et sans relancer l'optimisation.

Les exports PNG facilitent la consultation, tandis que les PDF et SVG
préservent des traits et textes vectoriels pour la rédaction. Les figures
du présent rapport utilisent des PDF du projet, régénérés depuis la synthèse
archivée lorsque nécessaire, sans nouvelle optimisation. Le schéma
de couverture reste l'image fournie par l'auteur, avec son crédit ResearchGate.

\section{Contrôles et tests}
La suite comprend 101 tests. Elle couvre notamment les paramètres invalides,
la lecture des coordonnées, les coûts, les contraintes des solutions,
les méthodes, l'interface, les exports et la disposition des figures.
Des cas limites vérifient les coordonnées dupliquées, un objectif nul,
$p=n$ et les arrêts sans échange possible.

Les contrôles les plus utiles sont ceux qui pourraient détecter un résultat
faux malgré un programme qui s'exécute : comparaison de petites PLNE à une
énumération indépendante, vérification de 2-opt sur l'arête de fermeture,
distinction entre solution entière sous budget et optimum prouvé, rejet
d'un coût sauvegardé incohérent, et relecture des solutions illustrées.
Les campagnes sont testées avec plusieurs poids pour éviter de fusionner
des cas différents dans les synthèses.

Les scripts d'archive et de nettoyage sont aussi testés. Ils doivent préserver
le dépôt Git, l'environnement local et les résultats actuels, et produire
une archive sans dossier enveloppe superflu. Ces vérifications concernent
l'usage du projet autant que ses algorithmes.

Pour la rédaction, les 172 solutions de la campagne ont été relues et leurs
coûts recalculés indépendamment à partir des coordonnées, des cycles et
des affectations. Les moyennes des tableaux ont été rapprochées des coûts
unitaires. Le jeu d'essais de 101 tests a été relancé avec le contrôle
d'export renforcé ; il se termine sans échec dans l'environnement de revue.
Cela ne garantit pas l'absence de tout défaut ni la compatibilité avec
toutes les versions possibles des dépendances.

\section{Un défaut d'export identifié pendant la revue}
Le fichier \code{completion.json} annonce 192 enregistrements et la synthèse
contient les 96 groupes attendus. Le CSV détaillé original n'en contient que
191 : il manque le cas \code{rd100}, $p=10$, PLNE, ignoré par configuration.
Le programme écrivait le CSV après les exécutions calculées, mais pas après
un dernier cas ignoré. Il s'agit donc d'une écriture finale manquante,
sans perte de solution calculée.

Le rapport conserve l'export original et produit une copie réconciliée distincte.
La ligne manquante est reconstituée comme cas ignoré, à partir de la
configuration, de la synthèse et du bilan de complétion. Aucun coût ni
temps de résolution n'est inventé pour ce cas. Les sources incluent le
correctif qui écrit tous les enregistrements à la fin de la campagne,
ainsi qu'une assertion supplémentaire dans le test existant des cas ignorés.
Ce correctif est joint pour intégration au dépôt ; il ne fait pas partie
des trois commits publiés mentionnés au début du rapport.

\section{Limites expérimentales et améliorations prioritaires}
Trois graines donnent un premier aperçu de la variabilité, mais restent peu
nombreuses pour quantifier finement la robustesse des recherches. Les petites
durées sont mesurées une fois par exécution, sans protocole dédié de répétition
ou d'isolation des ressources. La grille constitue une seule référence
constructive : un autre départ pourrait modifier les gains.

La priorité expérimentale serait de comparer les recherches à même nombre
d'évaluations, puis à même budget de temps, avec davantage de graines.
Il serait ensuite utile de faire varier $\alpha$ et $p$ de façon plus large,
d'essayer plusieurs initialisations et de mesurer séparément l'accès,
l'anneau et la distribution des distances individuelles. Une extension de
la PLNE pourrait étudier des relaxations plus fortes ou une formulation
avec séparation de coupes, en mesurant leur apport au-delà du seul temps.

Sur le plan algorithmique, la reconstruction complète de chaque voisin
pourrait être remplacée en partie par des mises à jour incrémentales des
affectations et des modifications locales du cycle. Ce gain de vitesse
demanderait des contrôles supplémentaires pour conserver les mêmes coûts
et ne pas introduire de dépendance cachée à l'état précédent.

Enfin, une application au transport demanderait des distances de réseau,
des demandes, des capacités, des contraintes d'implantation et une prise
en compte de l'équité. Les instances TSPLIB et les unités de coordonnées
ne permettent pas de conclure qu'un tracé proposé serait faisable ou
préférable pour une agglomération réelle.
\bilan{La traçabilité permet de vérifier les résultats et de corriger un
défaut d'export sans modifier les calculs. Les prochaines avancées utiles
portent sur l'équité des budgets, la robustesse des expériences et le
réalisme des données de transport.}
